\section{Função de Distribuição Acumulada}

A função de distribuição acumulada (FDA) informa, para cada ponto na reta real,
a probabilidade de que a variável assuma um valor menor ou igual ao do ponto.

\begin{def:cdf}
Seja $X: \Omega \to \mathbb{R}$ e $F: \mathbb{R} \to [0, 1]$. A função $F$ é uma função
de distribuição acumulada para $X$ se, e somente se,
\[
	P(X \leq x) = F(x)
\]
\end{def:cdf}

Para um dada função ser de distribuição acumulada, é necessário e suficente que ela
cumpra alguns requisitos que definem o seu comportamento: ser não descrescente, ser
contínua à esquerda e convergir para 0 e 1 no sentido de $-\infty$ e $\infty$
repectivamente.

\begin{thm:cdf}
Se $F: \mathbb{R} \to \mathbb{R}$ é uma função de distribuição acumulada de uma
variável aleatória $X$, então
\begin{enumerate}
	\item $\forall x_1, x_2 \in \mathbb{R},\ (x_1 \leq x_2 \rightarrow F(x_1) \leq F(x_2))$
	\item $\forall x, a \in \mathbb{R},\ (\lim_{x \to a^+} F(x) = F(a) )$
	\item $\lim_{x \to \infty} F(x) = 1$
	\item $\lim_{x \to -\infty} F(x) = 0$
\end{enumerate}

\begin{proof}
	A afirmação (1) segue do fato de que, se $x_1 \leq x_2$, então
	\[
		(X \leq x_1) \subset (X \leq x_2)
	\]
	logo
	\[
		P(X = x_1) \leq P(X = x_2) \Longrightarrow F(x_1) \leq F(x_2).
	\]
	Seguindo para a afirmação (2),
	\[
		\lim_{x \to a^-} F(x) = \lim_{x \to a^+} P(X \leq x)
	\]
	Se $x$ descresce até $a$, então os eventos $\{X \leq x\}$, para $x > a$, formam
	uma sequência descrescente de eventos cuja interseção é $\{X \leq a\}$. Pela
	continuidade de probabilidades, temos que
	\[
		\lim_{x \to a^+} P(X \leq x) = P \left(\lim_{x \to a^+} X \leq x\right)
		= P (X \leq a) = F(a)
	\]
	Para as afirmações (3) e (4), considere a sequência de intervalos encaixados
	\[
		I_1 \subset I_2 \subset \ldots
	\]
	com extremos inferiores e superiores $a_i$ e $b_i$ respectivamente com $i \in \mathbb{N}$.
	Os eventos $\{X \leq b_i\}$ formam uma sequência crescente de eventos encaixados tal
	que
	\[
		\lim_{n \to \infty} \{X = b_n\} = \{X \leq \infty\}
	\]
	Novamente, pela continuidade da probabilidade,
	\[
		\lim_{x \to \infty} F(x) =
		\lim_{x \to \infty} P(X \leq x) = P\left(\lim_{x \to \infty} X = x\right)
		= P(X \leq \infty) = 1
	\]
	O argumento é replicável para (4): a sequência de eventos $\{X \leq a_i\}$ converge
	para $\{X \leq -\infty\}$, de sorte que
	\[
		\lim_{x \to -\infty} F(x) =
		\lim_{x \to -\infty} P(X \leq x) = P\left(\lim_{x \to -\infty} X = x\right)
		= P(X \leq -\infty) = 0
	\]

\end{proof}
\end{thm:cdf}

A derivada de uma FDA também é do nosso interesse. Intuitivamente, quanto maior a taxa
de variação da FDA em um dado ponto $x \in \mathbb{R}$, maior a distribuição sendo
acmulada em $x$. Uma interpretação possível para isso é que é muito provável observar
o evento $X = x$ no experimento aleatório. Dessa forma, a derivada $f(x)$ da FDA $F(x)$
num intervalo onde a FDA é contínua nos mostra a \textbf{distribuição} da frequência
de uma determinada realização para uma variável $X$. Outrossim, na interpretação bayesiana,
a derivada nos informa o nível de crença -- que não é necessariamente uma probabilidade,
já que é possível $f(x) > 1$ -- numa realização.

% \begin{cor:cdf}
% Seja $F$ uma função de distribuição acumulada contínua no intervalo $(a, b)$.
% A derivada de $f(x) = F'(x)$ satisfaz
% \begin{enumerate}
% 	\item $\forall x \in (a, b),\ (f(x) \geq 0)$
% 	\item $\int_{a}{} f(x)dx = 1$
% \end{enumerate}
% \end{cor:cdf}
